\documentclass[11pt,a4paper]{article}
\usepackage[utf8]{inputenc}
\usepackage[T1]{fontenc}
\usepackage{amsmath,amssymb}
\usepackage{booktabs}
\usepackage{hyperref}
\usepackage{geometry}
\geometry{margin=2.5cm}
\usepackage{enumitem}

\title{\textbf{0\_data\_processing} \\ Experimental Dataset Construction Pipeline}
\author{AzONetV2}
\date{}

\begin{document}
\maketitle

\section{Overview}

This folder contains the pipeline to construct the experimental dataset by merging transmittance spectra ({\tt .txt}) with thickness measurements ({\tt .csv}), correcting thickness values where needed, and selecting the best-quality samples. The final output is a cleaned dataset of 145 AZO (Aluminium-doped Zinc Oxide) thin-film samples.

\section{Notebooks}

\subsection{Notebook 0 --- 101 Samples}
\textbf{File:} {\tt 0\_DataFrame\_Muestras\_Laboratorio.ipynb}

\textbf{Input:}
\begin{itemize}
    \item Thickness: {\tt ../../experimental\_samples/Processed\_Data/Thickness-Order-100-1600-3.csv} (101 samples, columns: Nombre, Nombre-Arc, Espesor, Error, Error Porcentual)
    \item Spectra: {\tt ../../experimental\_samples/Spectrum\_Order/*.txt}
\end{itemize}

\textbf{Process:}
\begin{itemize}
    \item Reads 101 samples with pre-corrected thickness values --- no linear transformation is applied.
    \item Uppercases Nombre-Arc for consistency.
    \item Generates histograms (bins of 3, 5, and 10), KDE density estimation, and error distribution plots.
    \item Verifies cross-matching: all 101 samples in the CSV match spectra in the folder (both unique sets are empty).
    \item Merges thickness + spectra on Nombre-Arc via inner join.
\end{itemize}

\textbf{Output:} {\tt ../../results/dataframe\_spectrum\_thickness\_101.pkl}

\subsection{Notebook 1 --- 15 Samples}
\textbf{File:} {\tt 1\_DataFrame\_Muestras\_Laboratorio\_15.ipynb}

\textbf{Input:}
\begin{itemize}
    \item Thickness: {\tt ../../experimental\_samples/Processed\_Data/Example\_Sample\_IA\_2\_F\_30NEW.csv} (65 rows, only 15 match spectra)
    \item Spectra: {\tt ../../experimental\_samples/Muestra\_ANZO67-ANZO69/*.txt}
\end{itemize}

\textbf{Process:}
\begin{itemize}
    \item Same structure as Notebook 0, but thickness requires a \textbf{linear correction}:
    \begin{equation*}
        E_{\text{corrected}} = \frac{E_{\text{raw}} - 17.897827}{1.1521356}
    \end{equation*}
    \item 50 CSV entries do not have corresponding spectra and are excluded by the inner join.
\end{itemize}

\textbf{Output:} {\tt ../../results/dataframe\_spectrum\_thickness\_15.pkl}

\subsection{Notebook 2 --- 45 Samples}
\textbf{File:} {\tt 2\_DataFrame\_Muestras\_Laboratorio\_45.ipynb}

\textbf{Input:}
\begin{itemize}
    \item Thickness: {\tt ../../experimental\_samples/Processed\_Data/Sample\_Test\_IA\_1\_New.csv} (45 rows)
    \item Spectra: {\tt ../../experimental\_samples/Sample\_test/*.txt} (51 files; 6 spectra without CSV match are excluded)
\end{itemize}

\textbf{Process:} Nearly identical to Notebook 1. Applies the \textbf{same linear correction} as Notebook 1.

\textbf{Output:} {\tt ../../results/dataframe\_spectrum\_thickness\_45.pkl}

\subsection{Notebook 3 --- Unify All Samples}
\textbf{File:} {\tt 3\_Unified\_101\_45\_15.ipynb}

\textbf{Input:} The three {\tt .pkl} files from Notebooks 0, 1, and 2.

\textbf{Process:}
\begin{itemize}
    \item Concatenates all dataframes into a single dataframe of \textbf{161 samples} (101 + 45 + 15).
    \item Adds a Nombre column to the 45-sample and 15-sample dataframes (where Nombre was the CSV index rather than a column) for schema consistency.
\end{itemize}

\textbf{Output:} {\tt ../../results/dataframe\_spectrum\_thickness\_161.pkl}

\subsection{Notebook 4 --- Figures}
\textbf{File:} {\tt 4\_Figures.ipynb}

\textbf{Input:}
\begin{itemize}
    \item {\tt ../../results/dataframe\_spectrum\_thickness\_161.pkl}
    \item Precomputed index array: {\tt ./indexes/outlier\_indices\_16.npy} (16 outlier indices)
\end{itemize}

\textbf{Process:}
\begin{itemize}
    \item Loads the 161-sample dataframe and drops the 16 outliers identified during manual inspection, yielding the final \textbf{145 samples}.
    \item Computes the mean percentage error: $\bar{\varepsilon} \approx 9.61\%$.
    \item Generates and saves 5 figures to {\tt ./images/}:
    \begin{itemize}
        \item Thickness histograms at 200 nm, 100 nm, and 50 nm bin widths (bin edges from 100 to 1500 nm).
        \item Percentage error histogram (1\% bins from 0 to 40\%).
        \item KDE density plot of the thickness distribution.
    \end{itemize}
\end{itemize}

\textbf{Output:}
\begin{itemize}
    \item {\tt ../../results/dataframe\_spectrum\_thickness\_145\_final.pkl} (the final, cleaned dataset)
    \item {\tt ./images/bin\_200.png}, {\tt bin\_100.png}, {\tt bin\_50.png}, {\tt error1.png}, {\tt kde.png}
\end{itemize}

\section{Key Differences Between Batches}

\begin{table}[htbp]
\centering
\begin{tabular}{lcccc}
\toprule
\textbf{Notebook} & \textbf{Samples} & \textbf{Thickness Correction} & \textbf{Spectra Source} \\
\midrule
0 & 101 & None (pre-corrected) & Spectrum\_Order/ \\
1 & 15  & Linear: $(x - 17.90)/1.15$ & Muestra\_ANZO67-ANZO69/ \\
2 & 45  & Linear: $(x - 17.90)/1.15$ & Sample\_test/ \\
\bottomrule
\end{tabular}
\end{table}

\begin{itemize}
    \item Notebooks 1 and 2 came from a different measurement batch, requiring the same thickness calibration.
    \item Notebooks 0--2 produce independent dataframes; Notebook 3 merges them.
    \item Notebook 4 drops 16 manually-identified outlier samples in a single step (161 $\rightarrow$ 145), then generates figures.
\end{itemize}

\section{Final Dataset}

The final dataset ({\tt dataframe\_spectrum\_thickness\_145\_final.pkl}) contains \textbf{145 samples}, each with the following columns:

\begin{itemize}
    \item \textbf{Nombre} \& \textbf{Nombre-Arc}: sample identifiers.
    \item \textbf{Espesor}: corrected thickness (nm).
    \item \textbf{Error} \& \textbf{Error Porcentual}: thickness measurement error (nm) and relative error (\%).
    \item \textbf{Espectro}: transmittance spectrum (911 wavelength points, 190--1100 nm).
\end{itemize}

\section{Dependencies}
\begin{itemize}
    \item {\tt pandas}, {\tt numpy}, {\tt matplotlib}, {\tt seaborn}
    \item {\tt pathlib}, {\tt re}
\end{itemize}

\end{document}
